\chapter{The Normal Distribution}

The previous chapter built distributions for specialized settings: waiting times,
proportions, and flat intervals. But when many small, independent influences add
together, a single shape appears across science: the smooth, symmetric bell curve.

It is the curve statisticians sketch on napkins, the shape of repeated measurement
errors, and the distribution of sample averages. It carries two names: the
\emph{normal} distribution, and the \emph{Gaussian} (after Carl Friedrich Gauss).
This chapter explores that curve: where the formula $e^{-x^2/2}$ comes from, why
$\sqrt{2\pi}$ sits out front, how $\mu$ and $\sigma$ control its center and width,
and why adding normals produces another normal.

\chapterline

\section{A Curve Born From Errors}
\label{ch12:errors}

Begin with a story, because the normal distribution was discovered through a very
human problem: \emph{nobody can measure anything exactly}.

Imagine an eighteenth-century astronomer trying to pin down the position of a
star. He looks through his telescope and writes down an angle. Then he looks
again and writes down a slightly different angle. Again -- different again. His
hand trembles, the air shimmers, his eye tires. Each measurement is the true
position plus a little \emph{error}, and the errors are different every time.

What can he say about these errors? A few things feel obvious once you say them
out loud:

\begin{itemize}
  \item Small errors are more common than large errors. A measurement off by a
  hair happens constantly; a measurement off by a mile almost never.
  \item Errors are symmetric: overshooting by $0.1$ is just as likely as
  undershooting by $0.1$. There is no built-in bias toward ``too big'' or ``too
  small'' if the instrument is unbiased.
  \item Errors pile up smoothly. There is no sharp threshold where an error
  suddenly becomes impossible; the chance simply decays as you move further out.
\end{itemize}

\begin{intuitionbox}
The astronomer's instinct -- ``small errors common, big errors rare, symmetric
about the truth, fading smoothly'' -- is not yet a formula. But it is already a
\emph{shape}: a single hump centered on the true value, tallest in the middle,
falling away on both sides, never quite touching zero. Every detail of the normal
density is just this picture made precise.
\end{intuitionbox}

The genius of Gauss (and before him, in pieces, of de Moivre and Laplace) was to
turn that intuition into an exact curve by adding one more natural-sounding
assumption -- and then asking what function is \emph{forced} on you. We will see
the resulting form in Section~\ref{ch12:standard}. But first, let us simply meet
the curve.

\begin{historybox}
The bell curve was not conceived in an armchair; it was forged across two
centuries of practical frustration with gambling odds, lost planets, and noisy
telescopes.

In 1733, Abraham de Moivre (1667--1754), a French Huguenot mathematician living in
London, was calculating probabilities for high-stakes gamblers. He wanted to
find the chance of obtaining various numbers of heads across hundreds of coin
flips -- a problem governed by the \emph{binomial distribution} (the formula for
the probability of getting $k$ successes in $n$ independent yes/no trials).
Calculating terms like $\binom{1000}{500}$ by hand meant multiplying thousands of
integers, an impossible chore before computers. Armed with early calculus and
James Stirling's approximation for factorials (the product of all integers up to
$n$), de Moivre discovered that as the number of coin tosses grows, the discrete
staircase of binomial probabilities hugs a continuous exponential curve:
$e^{-c x^2}$. It was the very first recorded sighting of the bell curve.

Seven decades later, in 1801, the Italian astronomer Giuseppe Piazzi spotted a
faint new celestial body, the dwarf planet Ceres. But after tracking it for only
$41$ days, Ceres slipped behind the glare of the Sun and was lost. Astronomers
across Europe tried to predict where it would reappear, but their calculations
conflicted wildly because every observatory's telescope recorded slightly
different, error-prone angles. Carl Friedrich Gauss (1777--1855), then 24 years
old, attacked the problem with a bold philosophical question: if an observer
makes repeated, noisy measurements of an unknown physical quantity, what
mathematical shape must the errors have so that the ordinary arithmetic average
(the sum divided by the count) is the most probable true value? Working backward
using calculus, Gauss proved that this single postulate uniquely forces errors to
follow the bell-shaped density $e^{-x^2/(2\sigma^2)}$. Combining this error law
with his \emph{method of least squares} -- a technique that fits a curve by
minimizing the sum of squared vertical gaps between the data points and the
model -- Gauss computed Ceres's orbit with extraordinary precision. When
astronomers aimed their telescopes where Gauss instructed in December 1801, Ceres
reappeared right on target.

A bitter controversy erupted almost immediately. When Gauss published his error
law in 1809, he claimed in passing that he had been using least squares since
1795. This infuriated the French mathematician Adrien-Marie Legendre, who had
published the method first in 1805 without a probability derivation, and felt
Gauss had stolen his priority. The theoretical debate deepened in 1810 when
Pierre-Simon Laplace (1749--1827) showed that Gauss's founding assumption -- that
the arithmetic average is inherently the best estimate -- was actually
unnecessary. Laplace proved via an early version of the \emph{Central Limit
Theorem} that the bell curve emerges automatically whenever an observed error is
the sum of many small, independent disturbances, regardless of what the
individual disturbance curves look like.

A second controversy arose over the name. In the late nineteenth century, Francis
Galton, Charles Sanders Peirce, and Karl Pearson popularized the term ``normal''
distribution (using ``normal'' in its geometric sense of perpendicular,
orthogonal, or standard). Pearson later regretted the term, observing that
anointing one curve as ``normal'' misled generations of scientists into believing
that any data with non-bell shapes was ``abnormal'' or defective.

What did this unlock? It provided the mathematical engine of the scientific
revolution: experimental error analysis, regression modeling in economics,
polling theory, and statistical mechanics (where James Clerk Maxwell used it in
1860 to describe the distribution of molecular speeds in a gas).

What is still open? In modern high-dimensional statistics and network science,
mathematicians actively study the exact rates of \emph{convergence} -- how many
random terms you must add before the sum genuinely resembles a bell, using tools
such as Stein's method for weakly dependent variables. Meanwhile, researchers
continue to map the precise boundaries where \emph{heavy-tailed} phenomena
(processes prone to sudden, catastrophic shocks, such as financial crashes or
power grid collapses) destroy Gaussian universality and demand new mathematics.
\end{historybox}

\section{The Standard Normal: The Simplest Bell}
\label{ch12:standard}

The cleanest version of the curve has no free knobs at all. It is centered at
zero, has its width fixed at one, and is called the \textbf{standard normal}. Its
density -- the function whose area under pieces gives probability -- is

\begin{formalbox}
\textbf{Standard normal density.} A random variable $Z$ is \emph{standard normal},
written $Z \sim \Normal(0,1)$, if it has density
\[
  \varphi(z) \;=\; \frac{1}{\sqrt{2\pi}}\, e^{-z^2/2}, \qquad z \in \R.
\]
We write $\varphi$ (lowercase phi) for this density and $\Phi$ (uppercase phi)
for its cumulative distribution function, $\Phi(z) = \P(Z \le z) = \int_{-\infty}^{z}\varphi(t)\,\dd t$.
The cumulative distribution function $\Phi(z)$ simply measures the total
probability accumulated to the left of $z$ -- the total shaded area under the
bell from $-\infty$ up to the threshold $z$.
\end{formalbox}

Let us read this formula slowly, because every symbol is doing a job.

\begin{itemize}
  \item The $e^{-z^2/2}$ is the \emph{heart}. The $z^2$ makes it symmetric (since
  $(-z)^2 = z^2$, the curve looks the same on both sides). The minus sign in the
  exponent makes it \emph{decrease} as you move away from zero. The squaring makes
  it decrease \emph{fast} -- faster than any straight line, faster even than a
  simple $e^{-|z|}$ -- so the tails are thin.
  \item The $\tfrac{1}{2}$ inside the exponent is a choice of scale that, as we
  will see, makes the variance come out to exactly $1$. Without it the curve would
  still be a bell, just a different width.
  \item The $\tfrac{1}{\sqrt{2\pi}}$ out front is a \emph{normalizing constant}.
  It is there for exactly one reason: to make the total area under the curve equal
  $1$, as every density must. It carries no shape information at all.
\end{itemize}

\begin{figure}[H]
\centering
\begin{tikzpicture}
\begin{axis}[
  width=12cm, height=6.5cm,
  domain=-4:4, samples=160,
  xlabel={$z$},
  ymin=0, ymax=0.46,
  xtick={-3,-2,-1,0,1,2,3},
  ytick=\empty,
  axis x line=bottom, axis y line=none,
  enlargelimits=false,
  title style={font=\small},
  every axis plot/.append style={thick},
]
\addplot[name path=curve, densityred, very thick]
  {1/sqrt(2*pi)*exp(-x^2/2)};
\path[name path=axis] (axis cs:-4,0) -- (axis cs:4,0);
\addplot[distblue!25] fill between[of=curve and axis, soft clip={domain=-1:1}];
\addplot[distblue!45] fill between[of=curve and axis, soft clip={domain=-2:-1}];
\addplot[distblue!45] fill between[of=curve and axis, soft clip={domain=1:2}];
\draw[dashed, black!40] (axis cs:0,0) -- (axis cs:0,0.3989);
\node[above=2pt] at (axis cs:0,0.3989) {\small $\varphi(z)$};
\node[fill=distblue!25, rounded corners=2pt, inner sep=2pt] at (axis cs:0,0.16) {\small $68\%$};
\node at (axis cs:1.5,0.06) {\small $13.6\%$};
\node at (axis cs:-1.5,0.06) {\small $13.6\%$};
\end{axis}
\end{tikzpicture}
\caption{The standard normal density $\varphi(z) = \frac{1}{\sqrt{2\pi}}e^{-z^2/2}$.
The shaded regions are the chance of landing within one and two standard
deviations of the center. The peak height is $1/\sqrt{2\pi}\approx 0.399$.}
\label{ch12:fig-standard}
\end{figure}

Notice in Figure~\ref{ch12:fig-standard} that the curve is tallest at $z=0$,
where its height is $1/\sqrt{2\pi}\approx 0.399$, and that it has already become
visually flat by $z = \pm 3$. By $z = \pm 4$ it is essentially glued to the
horizontal axis, though in truth it never reaches it -- the curve is positive for
\emph{every} real number, no matter how far out. There is, in principle, a
nonzero chance of an arbitrarily large error. It is just astronomically small.

\begin{walkbox}
The peak of $\varphi$ is about $0.399$. A narrow normal, say $\Normal(0, 0.1^2)$,
has a peak height of $1/(0.1\sqrt{2\pi})\approx 3.99$ -- a density value well above
$1$. But probabilities can never exceed $1$. How can a density take a value bigger
than any probability, and yet areas under it still come out as proper
probabilities? Pause and recall the difference between a density and a
probability. (Answer: a density is not a probability; it is probability \emph{per
unit length}. You multiply by a width to get probability, and the widths near a
tall narrow peak are tiny. The total area, not the peak height, is what must
equal $1$.)
\end{walkbox}

\subsection{Why Radial Symmetry Forces the Bell: The Herschel--Maxwell Dartboard}
\label{ch12:dartboard}

Where does that specific functional form $e^{-z^2/2}$ come from? Why not $e^{-|z|}$, or $1/(1+z^2)$, or some other symmetric decay?
In 1850 the astronomer John Herschel (and a decade later, the physicist James Clerk Maxwell in his study of gas molecule velocities) asked what shape of error distribution in two dimensions is geometrically natural. They imposed only two physical symmetries:

\begin{enumerate}
  \item \textbf{Orthogonal independence:} The horizontal deviation $X$ and the vertical deviation $Y$ of a dart thrown at the bullseye $(0,0)$ do not inform one another. The errors along perpendicular axes are independent, so their joint probability density must factor into separate 1D densities:
  \[
    f(x,y) = g(x)\,g(y).
  \]
  \item \textbf{Rotational symmetry:} The board has no preferred direction. Rotating the dartboard through any angle leaves the probability distribution unchanged. The joint density must depend strictly on the Euclidean distance from the origin, $r^2 = x^2 + y^2$:
  \[
    f(x,y) = \psi(x^2 + y^2).
  \]
\end{enumerate}

Now set these two expressions equal:
\[
  g(x)\,g(y) = \psi(x^2 + y^2).
\]
Setting $y=0$ gives $g(x)\,g(0) = \psi(x^2)$, which reveals that $g(x) = \psi(x^2)/g(0)$.
Substitute this back in:
\[
  \frac{\psi(x^2)}{g(0)} \cdot \frac{\psi(y^2)}{g(0)} = \psi(x^2 + y^2).
\]
Let $u = x^2 \ge 0$, $v = y^2 \ge 0$, and define $h(u) = \psi(u)/g(0)^2$. Then
\[
  h(u)\,h(v) = h(u+v).
\]
This is Cauchy's functional equation: a relationship where adding inputs
multiplies outputs ($h(u+v) = h(u)h(v)$). The only continuous functions on
$[0,\infty)$ that turn addition into multiplication and decay as distance
increases are decreasing exponentials:
\[
  h(u) = e^{-c u} \qquad (c > 0).
\]
Recalling $u = x^2$, the one-dimensional density $g(x)$ must satisfy
\[
  g(x) \;\propto\; e^{-c x^2}.
\]
Setting $c = \frac{1}{2\sigma^2}$ produces the Gaussian factor $e^{-x^2/(2\sigma^2)}$.
The bell curve is not a mathematical contrivance chosen for convenience. It is the \emph{unique} probability distribution whose multidimensional spread possesses both perpendicular independence and circular rotational symmetry (Figure~\ref{ch12:fig-dartboard}).

\begin{figure}[H]
\centering
\begin{tikzpicture}[scale=1.0]
% Coordinate axes
\draw[->, thick, black!70] (-3.4,0) -- (3.8,0) node[right, font=\footnotesize]{$x$};
\draw[->, thick, black!70] (0,-3.4) -- (0,3.6) node[above, font=\footnotesize]{$y$};

% Concentric rings of the dartboard with soft semantic fills
\fill[sourceblue!10] (0,0) circle (2.8);
\fill[sourceblue!25] (0,0) circle (1.9);
\fill[sourceblue!45] (0,0) circle (1.0);
\draw[thick, sourceblue!70!black] (0,0) circle (2.8);
\draw[thick, sourceblue!80!black] (0,0) circle (1.9);
\draw[thick, sourceblue!90!black] (0,0) circle (1.0);

% Center bullseye
\fill[densityred] (0,0) circle (3pt);
\node[below left, font=\scriptsize, black!80] at (0,0) {$(0,0)$};

% Circular level set annotation
\draw[->, >=stealth, thick, quotientorange!90!black] (1.4,1.4) -- (2.3,2.3);
\node[right, font=\scriptsize, align=left, text=quotientorange!95!black] at (2.2,2.4)
  {\textbf{Circular level curves}\\$r^2 = x^2 + y^2 = \text{const}$};

% Horizontal Gaussian slice profile above
\begin{scope}[yshift=4.3cm]
  \draw[->, black!60] (-3,0) -- (3,0) node[right, font=\tiny]{$x$};
  \draw[->, black!60] (0,0) -- (0,1.2) node[above, font=\tiny]{$g(x)$};
  \draw[very thick, densityred, domain=-2.8:2.8, samples=80] plot (\x, {exp(-\x*\x/1.8)});
  \node[font=\scriptsize, densityred!90!black, above] at (0,1.65) {Horizontal slice $g(x) \propto e^{-c x^2}$};
\end{scope}

% Vertical Gaussian slice profile to the right
\begin{scope}[xshift=5.6cm]
  \draw[->, black!60] (0,-3) -- (0,3) node[above, font=\tiny]{$y$};
  \draw[->, black!60] (0,0) -- (1.2,0) node[right, font=\tiny]{$g(y)$};
  \draw[very thick, targetgreen!80!black, domain=-2.8:2.8, samples=80] plot ({exp(-\x*\x/1.8)}, \x);
  \node[font=\scriptsize, targetgreen!80!black, rotate=-90, above] at (1.85,0) {Vertical slice $g(y) \propto e^{-c y^2}$};
\end{scope}
\end{tikzpicture}
\caption{The Herschel--Maxwell Gaussian dartboard. Rotational symmetry requires concentric circular contours $\psi(x^2+y^2)$. Orthogonal independence requires the joint density to factor as $g(x)g(y)$. The only functional family satisfying both geometric symmetries is the Gaussian $e^{-c(x^2+y^2)} = e^{-cx^2}e^{-cy^2}$.}
\label{ch12:fig-dartboard}
\end{figure}

\subsection{Why the constant is \texorpdfstring{$1/\sqrt{2\pi}$}{1/sqrt(2 pi)}}
\label{ch12:gaussianintegral}

Here is the single most famous integral in all of mathematics, and the reason
$\pi$ -- a constant about \emph{circles} -- shows up in a curve about \emph{errors}:

\begin{theorem}[The Gaussian integral]
\label{ch12:thm-gaussint}
\[
  \int_{-\infty}^{\infty} e^{-z^2/2}\,\dd z \;=\; \sqrt{2\pi}.
\]
\end{theorem}

\begin{whybox}
The function $e^{-z^2/2}$ has no elementary antiderivative -- you cannot write its
integral in terms of polynomials, exponentials, logs, and trig functions. So the
usual ``find the antiderivative and plug in the limits'' is hopeless. The trick,
due to Poisson, is to compute the \emph{square} of the integral and pass to two
dimensions, where the rotational symmetry of $e^{-(x^2+y^2)/2}$ lets polar
coordinates rescue you.

Let $I = \int_{-\infty}^{\infty} e^{-z^2/2}\,\dd z$. Then
\[
  I^2 = \left(\int_{-\infty}^{\infty} e^{-x^2/2}\,\dd x\right)
        \left(\int_{-\infty}^{\infty} e^{-y^2/2}\,\dd y\right)
      = \int_{-\infty}^{\infty}\!\int_{-\infty}^{\infty} e^{-(x^2+y^2)/2}\,\dd x\,\dd y.
\]
Now switch to polar coordinates: $x^2 + y^2 = r^2$, and the area element
$\dd x\,\dd y$ becomes $r\,\dd r\,\dd\theta$. That extra factor of $r$ is
essential: it provides the derivative factor that makes direct substitution possible:
\[
  I^2 = \int_{0}^{2\pi}\!\int_{0}^{\infty} e^{-r^2/2}\, r\,\dd r\,\dd\theta
      = \int_{0}^{2\pi}\!\Big[ -e^{-r^2/2}\Big]_{0}^{\infty}\dd\theta
      = \int_{0}^{2\pi} 1\,\dd\theta = 2\pi.
\]
Therefore $I = \sqrt{2\pi}$, and dividing $e^{-z^2/2}$ by this number makes its
total area exactly $1$. That is the entire reason for the $\sqrt{2\pi}$.
\end{whybox}

\begin{connectionbox}
The substitution $u = r^2/2$ (so $\dd u = r\,\dd r$) is where that extra factor
of $r$ from the polar area element gets absorbed. If you have read the companion calculus
book, this is the cleanest possible illustration of why the Jacobian of the
polar change of variables matters: drop the $r$ and the integral is impossible;
keep it and it falls out in one line. The appearance of $\pi$ is geometry leaking
into analysis -- it comes from the $2\pi$ of a full rotation, the angular sweep of
the plane.
\end{connectionbox}

\begin{walkbox}
\textbf{Mental check: scaling the Gaussian integral.}
What is $\int_{-\infty}^{\infty} e^{-z^2}\,\dd z$? Notice there is no factor of
$1/2$ in the exponent. Pause and try to evaluate it in your head using
Theorem~\ref{ch12:thm-gaussint}. (Answer: Substitute $u = \sqrt{2}\,z$, so
$\dd z = \dd u/\sqrt{2}$. The integral becomes
$\frac{1}{\sqrt{2}}\int_{-\infty}^{\infty} e^{-u^2/2}\,\dd u = \frac{\sqrt{2\pi}}{\sqrt{2}} = \sqrt{\pi}$.
This value $\sqrt{\pi}$ explains why the error function $\operatorname{erf}$,
which integrates $e^{-t^2}$, is normalized by $2/\sqrt{\pi}$.)
\end{walkbox}

\section{Adding Two Knobs: Mean and Variance}
\label{ch12:general}

The standard normal is centered at $0$ with spread $1$. Real data are centered
wherever they happen to be centered and spread however they spread. So we give
the curve two adjustable knobs: a location $\mu$ and a scale $\sigma$.

\begin{formalbox}
\textbf{General normal density.} A random variable $X$ is \emph{normal} with mean
$\mu$ and variance $\sigma^2$, written $X \sim \Normal(\mu, \sigma^2)$, if it has
density
\[
  f(x) \;=\; \frac{1}{\sigma\sqrt{2\pi}}\,
              \exp\!\left(-\frac{(x-\mu)^2}{2\sigma^2}\right),
  \qquad x \in \R.
\]
Here $\mu \in \R$ is the \emph{mean} and $\sigma > 0$ is the \emph{standard
deviation}; $\sigma^2$ is the \emph{variance}.
\end{formalbox}

This is the same bell, repositioned and rescaled. Compare it term by term with
the standard normal and you can read off what each knob does:

\begin{itemize}
  \item Replacing $z$ by $(x-\mu)/\sigma$ \emph{slides} the curve so its peak sits
  at $x=\mu$ instead of $0$. The mean $\mu$ is the center of symmetry, the peak,
  the median, and (because the curve is symmetric) the average -- all at once.
  \item Dividing the deviation by $\sigma$ \emph{stretches} the horizontal axis.
  Large $\sigma$ spreads the curve wide and flat; small $\sigma$ squeezes it tall
  and narrow.
  \item The leading $\tfrac{1}{\sigma\sqrt{2\pi}}$ keeps the total area at $1$ even
  as you stretch: a wider curve must be shorter so the area is conserved. That is
  why a narrow normal is tall and a wide one is squat.
\end{itemize}

\begin{figure}[H]
\centering
\begin{tikzpicture}
\begin{axis}[
  width=12cm, height=6.5cm,
  domain=-8:10, samples=180,
  xlabel={$x$}, ylabel={density},
  ymin=0, ymax=0.45,
  axis lines=middle, enlargelimits=false,
  title style={font=\small},
  title={Three normal densities with varying $\mu$ and $\sigma$},
  legend pos=north east, legend cell align=left,
]
\addplot[densityred, very thick] {1/(1*sqrt(2*pi))*exp(-(x-0)^2/(2*1^2))};
\addlegendentry{$\Normal(0,1)$}
\addplot[distblue, very thick, dashed] {1/(2*sqrt(2*pi))*exp(-(x-0)^2/(2*2^2))};
\addlegendentry{$\Normal(0,4)$}
\addplot[targetgreen, very thick, dotted] {1/(1.2*sqrt(2*pi))*exp(-(x-4)^2/(2*1.2^2))};
\addlegendentry{$\Normal(4,1.44)$}
\end{axis}
\end{tikzpicture}
\caption{Three normals. Changing $\mu$ slides the bell left or right; changing
$\sigma$ makes it narrow-and-tall or wide-and-short. The area under each is
always exactly $1$.}
\label{ch12:fig-three}
\end{figure}

\subsection{The Geometric Meaning of $\sigma$: The Inflection Points}
\label{ch12:inflection}

Where does the standard deviation $\sigma$ live on the drawing of a bell curve?
With the uniform distribution, the spread was the width of the box. But the normal curve stretches infinitely in both directions, never quite reaching zero. What landmark on the curve corresponds to $\pm \sigma$?

The answer is given by calculus: $\mu \pm \sigma$ are the \emph{inflection points} of the density curve, where the curvature changes sign. Differentiating $f(x) = \frac{1}{\sigma\sqrt{2\pi}}\exp\left(-\frac{(x-\mu)^2}{2\sigma^2}\right)$ with respect to $x$:
\begin{align*}
  f'(x) &= -\frac{x-\mu}{\sigma^2}\,f(x), \\
  f''(x) &= -\frac{1}{\sigma^2}\,f(x) - \frac{x-\mu}{\sigma^2} f'(x)
          = \left(\frac{(x-\mu)^2 - \sigma^2}{\sigma^4}\right) f(x).
\end{align*}
Because $f(x) > 0$ everywhere, $f''(x) = 0$ if and only if
\[
  (x - \mu)^2 = \sigma^2 \quad\Longleftrightarrow\quad x = \mu \pm \sigma.
\]
This clean calculation gives us an immediate visual rule for reading $\sigma$:
\begin{itemize}
  \item In the central region $\abs{x - \mu} < \sigma$, $f''(x) < 0$: the curve is \textbf{concave down}, curving like the crest of a hill or an umbrella.
  \item At the shoulders $x = \mu \pm \sigma$, the slope is steepest: $|f'(\mu \pm \sigma)| = \frac{1}{\sigma^2 \sqrt{2\pi e}}$.
  \item In the outer flanks $\abs{x - \mu} > \sigma$, $f''(x) > 0$: the curve is \textbf{concave up}, bending outward like a trumpet to hug the horizontal axis.
\end{itemize}
The parameter $\sigma$ is the radius of the dome. When you see a bell curve, look for where the curve stops holding water like a bowl and begins flaring out: that transition occurs at exactly one standard deviation from the peak (Figure~\ref{ch12:fig-inflection}).

\begin{figure}[H]
\centering
\begin{tikzpicture}
\begin{axis}[
  width=12cm, height=6.2cm,
  domain=-4:4, samples=160,
  xlabel={$z = (x-\mu)/\sigma$},
  ymin=0, ymax=0.48,
  xtick={-3,-2,-1,0,1,2,3},
  xticklabels={$-3\sigma$,$-2\sigma$,$-\sigma$,$\mu$,$+\sigma$,$+2\sigma$,$+3\sigma$},
  axis x line=bottom, axis y line=none,
  ytick=\empty, enlargelimits=false,
  title style={font=\small},
  title={Inflection points and concavity regimes of the normal density}
]
\addplot[name path=curve, densityred, very thick] {1/sqrt(2*pi)*exp(-x^2/2)};

% Shaded central dome (concave down)
\path[name path=axis] (axis cs:-4,0) -- (axis cs:4,0);
\addplot[sourceblue!20] fill between[of=curve and axis, soft clip={domain=-1:1}];

% Inflection points markers
\draw[fill=distblue] (axis cs:1,0.242) circle (2.5pt);
\draw[fill=distblue] (axis cs:-1,0.242) circle (2.5pt);

% Tangent lines at inflection points
\addplot[thick, dashed, quotientorange, domain=0.4:1.7] {-0.242*(x-1) + 0.242};
\addplot[thick, dashed, quotientorange, domain=-1.7:-0.4] {0.242*(x+1) + 0.242};

% Labels in dedicated whitespace
\node[above=2pt, font=\footnotesize] at (axis cs:0, 0.3989) {$f(\mu)$};
\node[above, font=\footnotesize\bfseries, sourceblue!90!black, align=center, fill=sourceblue!20, rounded corners=2pt, inner sep=2pt] at (axis cs:0,0.06) {Concave Down ($f'' < 0$)\\[-2pt]\scriptsize (Dome / Cap)};
\node[above right, font=\footnotesize, distblue!90!black] at (axis cs:1.05,0.25) {Inflection Point ($+\sigma$)};
\node[above left, font=\footnotesize, distblue!90!black] at (axis cs:-1.05,0.25) {Inflection Point ($-\sigma$)};
\node[font=\footnotesize, black!70] at (axis cs:2.8,0.15) {Concave Up ($f'' > 0$)};
\node[font=\footnotesize, black!70] at (axis cs:-2.8,0.15) {Concave Up ($f'' > 0$)};
\end{axis}
\end{tikzpicture}
\caption{The inflection points of the normal curve occur at exactly $\mu \pm \sigma$. In the shaded region $|x-\mu|<\sigma$ the curve is concave down ($f''<0$). Beyond $\pm\sigma$ the curve turns concave up ($f''>0$), flattening out into the tails. The dashed orange lines show the tangents at the inflection points, where the descent is steepest.}
\label{ch12:fig-inflection}
\end{figure}

\begin{happeningbox}
Every normal distribution has the \emph{same} underlying bell curve. The two
parameters $\mu$ and $\sigma$ do not alter the intrinsic shape -- they only set
the center of the curve and how widely it spreads. This is why a single table or
a single function $\Phi$ is enough to answer any question about any normal: you
just translate and rescale. In mathematical terms, all normal distributions form
a single \emph{location-scale family} -- a family generated by shifting
(the location $\mu$) and scaling (the scale $\sigma$) a single base reference
density -- and the next subsection makes this transformation explicit.
\end{happeningbox}

\subsection{Standardizing: the universal translator}
\label{ch12:standardize}

The link between the general normal and the standard one is a single, enormously
useful operation: subtract the mean, divide by the standard deviation.

\begin{theorem}[Standardization]
\label{ch12:thm-standardize}
If $X \sim \Normal(\mu, \sigma^2)$ then
\[
  Z = \frac{X - \mu}{\sigma} \sim \Normal(0,1).
\]
Conversely, if $Z \sim \Normal(0,1)$ then $X = \mu + \sigma Z \sim \Normal(\mu,\sigma^2)$.
\end{theorem}

\begin{whybox}
This is a change of variables. Let $X \sim \Normal(\mu,\sigma^2)$ and set
$Z = (X-\mu)/\sigma$, so $X = \mu + \sigma Z$ and $\dd x = \sigma\,\dd z$. The
density of $Z$ is the density of $X$ at the point $x = \mu + \sigma z$, multiplied
by the stretch factor $|\dd x/\dd z| = \sigma$:
\[
  f_Z(z) = f_X(\mu + \sigma z)\cdot \sigma
         = \frac{1}{\sigma\sqrt{2\pi}}\exp\!\left(-\frac{(\mu+\sigma z - \mu)^2}{2\sigma^2}\right)\cdot \sigma
         = \frac{1}{\sqrt{2\pi}} e^{-z^2/2}.
\]
The $\sigma$ in the denominator cancels the $\sigma$ from the stretch, the
$(\mu + \sigma z - \mu)^2 = \sigma^2 z^2$ collapses, and out drops $\varphi(z)$
exactly. So $Z$ is standard normal.
\end{whybox}

Standardization is what lets you answer \emph{any} normal question with one table.
To find $\P(X \le x)$ for $X \sim \Normal(\mu,\sigma^2)$, you convert the
threshold $x$ into a number of standard deviations away from the mean -- a
\emph{$z$-score} (a dimensionless count of standard deviations above or below the mean) -- and look that up:
\[
  \P(X \le x) = \P\!\left(\frac{X-\mu}{\sigma} \le \frac{x-\mu}{\sigma}\right)
              = \Phi\!\left(\frac{x-\mu}{\sigma}\right).
\]

\begin{figure}[H]
\centering
\begin{tikzpicture}
\begin{axis}[
  name=topaxis,
  width=11.5cm, height=4.2cm,
  domain=200:800, samples=120,
  xlabel={$x$ (original exam score)}, ylabel={$f_X(x)$},
  ymin=0, ymax=0.005,
  xmin=150, xmax=850,
  xtick={200,350,500,650,800},
  xticklabels={$200$,$350$,$\mu=500$,$x=650$,$800$},
  x tick label style={font=\scriptsize},
  axis lines=left, enlargelimits=false,
  title style={font=\small},
  title={Original scale: $X \sim \Normal(500, 100^2)$},
]
\addplot[name path=curveA, distblue, very thick] {1/(100*sqrt(2*pi))*exp(-(x-500)^2/(2*10000))};
\path[name path=axisA] (axis cs:150,0) -- (axis cs:850,0);
\addplot[distblue!30] fill between[of=curveA and axisA, soft clip={domain=150:650}];
\draw[thick, dashed, densityred] (axis cs:650,0) -- (axis cs:650,0.0013);
\node[font=\footnotesize, distblue!90!black] at (axis cs:430,0.0016) {$\P(X \le 650) \approx 0.9332$};
\end{axis}

\begin{axis}[
  name=botaxis,
  at={($(topaxis.south)-(0,2.3cm)$)},
  anchor=north,
  width=11.5cm, height=4.2cm,
  domain=-3.5:3.5, samples=120,
  xlabel={$z$ ($z$-score in standard deviations)}, ylabel={$\varphi(z)$},
  ymin=0, ymax=0.48,
  xmin=-3.5, xmax=3.5,
  xtick={-3,-2,-1,0,1.5,3},
  xticklabels={$-3$,$-2$,$-1$,$0$,$z=1.5$,$3$},
  x tick label style={font=\scriptsize},
  axis lines=left, enlargelimits=false,
  title style={font=\small},
  title={Standardized scale: $Z = \frac{X - \mu}{\sigma} \sim \Normal(0,1)$},
]
\addplot[name path=curveB, targetgreen!80!black, very thick] {1/sqrt(2*pi)*exp(-x^2/2)};
\path[name path=axisB] (axis cs:-3.5,0) -- (axis cs:3.5,0);
\addplot[targetgreen!30] fill between[of=curveB and axisB, soft clip={domain=-3.5:1.5}];
\draw[thick, dashed, densityred] (axis cs:1.5,0) -- (axis cs:1.5,0.13);
\node[font=\footnotesize, targetgreen!90!black] at (axis cs:-0.7,0.15) {$\Phi(1.5) \approx 0.9332$};
\end{axis}

\draw[->, >=stealth, very thick, densityred] ($(topaxis.east)+(0.4cm,-0.3cm)$) -- ($(botaxis.east)+(0.4cm,0.3cm)$)
  node[midway, right=4pt, font=\footnotesize, align=left] {Standardize:\\$z = \dfrac{x - \mu}{\sigma}$\\[2pt]\scriptsize (Shift by $-\mu$,\\scale by $1/\sigma$)};
\end{tikzpicture}
\caption{Standardization transforms any normal curve into the standard normal template without altering probabilities. Shifting by $-\mu$ centers the curve at $0$; dividing by $\sigma$ rescales the width to $1$. The shaded area below $x = 650$ on $\Normal(500,100^2)$ matches the shaded area below $z = 1.5$ on $\Normal(0,1)$ exactly: $\P(X \le 650) = \Phi(1.5) \approx 0.9332$.}
\label{ch12:fig-standardize}
\end{figure}

\begin{example}[A test-score calculation]
\label{ch12:ex-zscore}
Scores on a standardized exam are modeled as $X \sim \Normal(500, 100^2)$. What
fraction of test-takers score above $650$?

The $z$-score of $650$ is
\[
  z = \frac{650 - 500}{100} = 1.5.
\]
So a score of $650$ is exactly $1.5$ standard deviations above the mean. Then
\[
  \P(X > 650) = \P(Z > 1.5) = 1 - \Phi(1.5) \approx 1 - 0.9332 = 0.0668.
\]
About $6.7\%$ of test-takers score above $650$ (illustrated in Figure~\ref{ch12:fig-standardize}). Notice we never needed a special
table for ``mean $500$, sd $100$'' -- standardizing reduced it to a single lookup
of $\Phi(1.5)$.
\end{example}

\begin{walkbox}
\textbf{Mental check: z-scores in your head.}
Adult systolic blood pressure follows roughly $\Normal(120, 15^2)$ mmHg.
Without looking at a table, what is the $z$-score of a reading of $150$ mmHg?
What about $105$ mmHg?
(Answer: For $150$, $z = (150 - 120)/15 = +2.0$: exactly two standard deviations above the average. For $105$, $z = (105 - 120)/15 = -1.0$: one standard deviation below. By the empirical rule, you know immediately that about $95\%$ of adults fall between $90$ and $150$ mmHg.)
\end{walkbox}

\section{Why $\mu$ and $\sigma^2$ Really Are the Mean and Variance}
\label{ch12:moments}

We have been \emph{calling} $\mu$ the mean and $\sigma^2$ the variance. We owe you
a proof that the names are accurate -- that if you compute $\E[X]$ and $\Var(X)$ from
the density, you actually get $\mu$ and $\sigma^2$.

\begin{proposition}
\label{ch12:prop-moments}
If $X \sim \Normal(\mu,\sigma^2)$ then $\E[X] = \mu$ and $\Var(X) = \sigma^2$.
\end{proposition}

\begin{whybox}
By standardization it is enough to handle $Z \sim \Normal(0,1)$ and then transfer
the result, because $X = \mu + \sigma Z$ gives $\E[X] = \mu + \sigma\E[Z]$ and
$\Var(X) = \sigma^2 \Var(Z)$ by linearity of expectation and the scaling rule for
variance.

\textbf{The mean of $Z$ is $0$.} The integrand $z\,\varphi(z)$ is an
\emph{odd} function: replacing $z$ by $-z$ flips its sign, because $z$ is odd and
$\varphi(z) = \varphi(-z)$ is even. The integral of an odd function over a
symmetric interval $(-\infty,\infty)$ is zero (the left half exactly cancels the
right half). Hence $\E[Z] = \int z\,\varphi(z)\,\dd z = 0$. So $\E[X] = \mu$.

\textbf{The variance of $Z$ is $1$.} Since $\E[Z]=0$, $\Var(Z) = \E[Z^2] =
\int_{-\infty}^{\infty} z^2 \varphi(z)\,\dd z$. Compute it by integration by
parts. Write $z^2 e^{-z^2/2} = z\cdot\big(z e^{-z^2/2}\big)$ and note that
$z\,e^{-z^2/2}$ is the derivative of $-e^{-z^2/2}$. With $u = z$ and
$\dd v = z e^{-z^2/2}\dd z$, so $v = -e^{-z^2/2}$:
\[
  \int_{-\infty}^{\infty} z^2 e^{-z^2/2}\,\dd z
   = \Big[-z\,e^{-z^2/2}\Big]_{-\infty}^{\infty}
     + \int_{-\infty}^{\infty} e^{-z^2/2}\,\dd z.
\]
The boundary term vanishes because as $z \to \pm\infty$, the quadratic
exponential decay in the denominator completely overpowers the linear factor:
$\lim_{z\to\pm\infty} \frac{-z}{e^{z^2/2}} = 0$ (which can be confirmed
directly by L'H\^opital's rule). The remaining integral is the Gaussian
integral $\sqrt{2\pi}$ from Theorem~\ref{ch12:thm-gaussint}. Dividing by the
normalizer $\sqrt{2\pi}$,
\[
  \Var(Z) = \frac{1}{\sqrt{2\pi}}\cdot \sqrt{2\pi} = 1.
\]
Therefore $\Var(X) = \sigma^2 \cdot 1 = \sigma^2$, and the parameter names are
vindicated.
\end{whybox}

\begin{trapbox}
A persistent source of confusion: $\Normal(\mu,\sigma^2)$ is parameterized by the
\emph{variance} $\sigma^2$, not the standard deviation $\sigma$. So $\Normal(0,9)$
has standard deviation $3$, not $9$. Some books and software (notably some R
functions) instead take $\sigma$ directly. Always check which convention you are
in before plugging numbers, or you will be off by a square.
\end{trapbox}

\section{The 68--95--99.7 Rule}
\label{ch12:empirical}

Because every normal is the same curve rescaled, the probability of landing within
a given number of standard deviations of the mean is \emph{always the same}, no
matter what $\mu$ and $\sigma$ are. Three of these numbers are worth memorizing
for life; they are sometimes called the \emph{empirical rule}.

\begin{focusbox}
For \emph{any} $X \sim \Normal(\mu, \sigma^2)$:
\begin{align*}
  \P(\mu - 1\sigma \le X \le \mu + 1\sigma) &\approx 0.6827 \quad (\text{about } 68\%),\\
  \P(\mu - 2\sigma \le X \le \mu + 2\sigma) &\approx 0.9545 \quad (\text{about } 95\%),\\
  \P(\mu - 3\sigma \le X \le \mu + 3\sigma) &\approx 0.9973 \quad (\text{about } 99.7\%).
\end{align*}
\end{focusbox}

These come straight from the standard normal CDF:
$\P(|Z| \le 1) = \Phi(1) - \Phi(-1) = 2\Phi(1) - 1 \approx 2(0.84134) - 1 \approx 0.6827$,
and likewise for $2$ and $3$.

\begin{figure}[H]
\centering
\begin{tikzpicture}
\begin{axis}[
  width=13cm, height=6cm,
  domain=-4:4, samples=160,
  xlabel={standard deviations from the mean},
  ylabel={$\varphi(z)$},
  ymin=0, ymax=0.45,
  xtick={-3,-2,-1,0,1,2,3},
  xticklabels={$\mu-3\sigma$,$\mu-2\sigma$,$\mu-\sigma$,$\mu$,$\mu+\sigma$,$\mu+2\sigma$,$\mu+3\sigma$},
  x tick label style={font=\footnotesize, rotate=30, anchor=east},
  axis lines=left, enlargelimits=false, clip=false,
  title style={font=\small},
  title={The empirical rule: probability mass in standard deviation bands}
]
\addplot[name path=curve, densityred, very thick] {1/sqrt(2*pi)*exp(-x^2/2)};
\path[name path=axis] (axis cs:-4,0) -- (axis cs:4,0);
\addplot[distblue!20] fill between[of=curve and axis, soft clip={domain=-3:-2}];
\addplot[distblue!20] fill between[of=curve and axis, soft clip={domain=2:3}];
\addplot[distblue!45] fill between[of=curve and axis, soft clip={domain=-2:-1}];
\addplot[distblue!45] fill between[of=curve and axis, soft clip={domain=1:2}];
\addplot[distblue!70] fill between[of=curve and axis, soft clip={domain=-1:1}];
\node[white] at (axis cs:0,0.15) {\footnotesize $68\%$};
\node at (axis cs:1.5,0.06) {\footnotesize $13.6\%$};
\node at (axis cs:-1.5,0.06) {\footnotesize $13.6\%$};
\node at (axis cs:2.5,0.045) {\footnotesize $2.1\%$};
\node at (axis cs:-2.5,0.045) {\footnotesize $2.1\%$};
\end{axis}
\end{tikzpicture}
\caption{The 68--95--99.7 rule. Each shell of width one standard deviation adds
less and less probability; almost everything lives within three deviations of the
mean.}
\label{ch12:fig-empirical}
\end{figure}

\begin{intuitionbox}
With those three numbers in mind, you can sanity-check claims about normal data
on the spot. ``IQ is $\Normal(100, 15^2)$, and
someone scored $145$'' -- that is three standard deviations above the mean, so
they are out past the $99.85$th percentile, rarer than $1$ in $700$. ``A part is
supposed to be $50$mm with sd $0.5$mm, and yours measures $52$mm'' -- four
standard deviations out, which for a normal process should essentially never
happen, so your machine is probably broken. You did both of those in your head.
\end{intuitionbox}

\begin{applicationbox}
\textbf{Six Sigma.} The manufacturing-quality movement named ``Six Sigma'' takes
its name directly from this rule. The goal is to make a process so tightly
controlled that the nearest specification limit is six standard deviations away
from the mean. For a perfectly centered normal process, being outside
$\pm 6\sigma$ has probability about $2 \times 10^{-9}$ -- two defects per billion.
(In practice the program allows for a $1.5\sigma$ drift of the mean over time,
which is why the famous quoted figure is $3.4$ defects per million rather than
the pure-normal two per billion. The arithmetic of those tail areas is exactly
what you are learning to do here.)
\end{applicationbox}

\section{The Tails: Reading the Standard Normal Table}
\label{ch12:table}

Because $\Phi$ has no closed form, for centuries people read its values off a
printed table. Understanding the \emph{idea} of the table is worth more than any
particular software call, because it forces you to think in terms of $z$-scores
and tail areas.

A standard normal table simply lists $\Phi(z) = \P(Z \le z)$ for a grid of
$z$-values. A small slice:

\begin{center}
\begin{tabular}{ccccccc}
\toprule
$z$ & $0.00$ & $0.50$ & $1.00$ & $1.50$ & $2.00$ & $2.50$\\
\midrule
$\Phi(z)$ & $0.5000$ & $0.6915$ & $0.8413$ & $0.9332$ & $0.9772$ & $0.9938$\\
\bottomrule
\end{tabular}
\end{center}

The table is really a picture of one curve: the cumulative distribution function
$\Phi$, which climbs from $0$ on the far left to $1$ on the far right, passing
through $\tfrac12$ at the symmetric center $z=0$.

\begin{figure}[H]
\centering
\begin{tikzpicture}
\begin{axis}[
  width=12cm, height=6cm,
  xlabel={$z$}, ylabel={$\Phi(z)=\P(Z\le z)$},
  xmin=-3.5, xmax=3.5, ymin=0, ymax=1.08,
  xtick={-3,-2,-1,0,1,2,3},
  ytick={0,0.25,0.5,0.75,1},
  axis lines=left,
  title style={font=\small},
  title={The standard normal cumulative distribution function $\Phi(z)$},
]
\addplot[distblue, very thick, smooth] coordinates {
  (-3.500,0.0002) (-3.250,0.0006) (-3.000,0.0013) (-2.750,0.0030)
  (-2.500,0.0062) (-2.250,0.0122) (-2.000,0.0228) (-1.750,0.0401)
  (-1.500,0.0668) (-1.250,0.1056) (-1.000,0.1587) (-0.750,0.2266)
  (-0.500,0.3085) (-0.250,0.4013) (0.000,0.5000) (0.250,0.5987)
  (0.500,0.6915) (0.750,0.7734) (1.000,0.8413) (1.250,0.8944)
  (1.500,0.9332) (1.750,0.9599) (2.000,0.9772) (2.250,0.9878)
  (2.500,0.9938) (2.750,0.9970) (3.000,0.9987) (3.250,0.9994) (3.500,0.9998)
};
\draw[gray, dashed] (axis cs:0,0) -- (axis cs:0,0.5) -- (axis cs:-3.5,0.5);
\addplot[only marks, mark=*, mark size=2pt, distblue] coordinates {(0,0.5)};
\node[anchor=south east, font=\footnotesize, fill=white, inner sep=1.5pt] at (axis cs:-0.1,0.52) {$\Phi(0)=0.5$};
\end{axis}
\end{tikzpicture}
\caption{The standard normal CDF $\Phi$. It is the running area under the bell of
Figure~\ref{ch12:fig-standard}: an S-shaped (sigmoid) curve rising from $0$ to
$1$, steepest where the density is tallest (at $z=0$) and flattening into the
thin tails. Every table value is one point on this curve.}
\label{ch12:fig-cdf}
\end{figure}

Everything else you compute from three symmetry-and-complement facts:

\begin{focusbox}
\begin{enumerate}
  \item \textbf{Total area:} $\P(Z \le z) + \P(Z > z) = 1$, so $\P(Z > z) = 1 - \Phi(z)$.
  \item \textbf{Symmetry:} $\Phi(-z) = 1 - \Phi(z)$, because the curve is a mirror
  image about $0$.
  \item \textbf{Intervals:} $\P(a \le Z \le b) = \Phi(b) - \Phi(a)$.
\end{enumerate}
\end{focusbox}

\begin{example}[Working the table both ways]
\label{ch12:ex-table}
Let $Z \sim \Normal(0,1)$.

\emph{(a) A right tail.} $\P(Z > 1.5) = 1 - \Phi(1.5) = 1 - 0.9332 = 0.0668$.

\emph{(b) A left tail by symmetry.} $\P(Z \le -1) = \Phi(-1) = 1 - \Phi(1) =
1 - 0.8413 = 0.1587$.

\emph{(c) A symmetric middle.} $\P(-2 \le Z \le 2) = \Phi(2) - \Phi(-2)
= 0.9772 - (1 - 0.9772) = 0.9544$ -- the $95\%$ of the empirical rule, recovered
from the table.

\emph{(d) Inverse lookup.} ``What cutoff $c$ has $5\%$ of the mass above it?'' We
need $\Phi(c) = 0.95$, and the table (read backwards) gives $c \approx 1.645$.
This $1.645$ is the famous one-sided $5\%$ critical value you will meet again in
hypothesis testing.
\end{example}

\begin{connectionbox}
Modern software replaces printed tables with two functions: a CDF call
(\texttt{Phi}, e.g.\ \texttt{scipy.stats.norm.cdf}) and its inverse, the
\emph{quantile function} or \emph{probit} (\texttt{scipy.stats.norm.ppf}). The
quantile function takes a target probability $p \in (0,1)$ and returns the
threshold value $z$ satisfying $\Phi(z) = p$; the inverse lookup in
Example~\ref{ch12:ex-table}(d) is exactly $\Phi^{-1}(0.95) \approx 1.645$. The
probit function also reappears as a link function in statistics (in probit
regression, a statistical model that predicts binary yes/no probabilities
through the cumulative normal curve) and as the mathematical engine of the
inverse-CDF method for simulating normal random variables on a computer.
\end{connectionbox}

\begin{remark}[The error function]
\label{ch12:erf}
You will sometimes see $\Phi$ written through the \emph{error function}
$\operatorname{erf}(x) = \frac{2}{\sqrt{\pi}}\int_0^x e^{-t^2}\,\dd t$. The
relationship is $\Phi(z) = \tfrac{1}{2}\big(1 + \operatorname{erf}(z/\sqrt{2})\big)$.
The two functions differ only by a change of variable: the factor $\sqrt{2}$
reflects the difference between the $e^{-t^2}$ and $e^{-z^2/2}$ scaling conventions.
\end{remark}

\section{Closure Properties: Normality Survives Linear Operations}
\label{ch12:closure}

Here is what makes the normal distribution not merely common but \emph{convenient}.
Two operations that destroy the shape of most distributions leave the normal
completely unharmed.

\subsection{Linear transformations stay normal}
\label{ch12:linear}

We already saw half of this in standardization. The full statement:

\begin{theorem}[Linear closure]
\label{ch12:thm-linear}
If $X \sim \Normal(\mu, \sigma^2)$ and $a \ne 0$, $b$ are constants, then
\[
  aX + b \sim \Normal\!\big(a\mu + b,\; a^2\sigma^2\big).
\]
\end{theorem}

The mean shifts and scales the way any mean does ($\E[aX+b] = a\mu + b$); the
variance scales by $a^2$ ($\Var(aX+b) = a^2\sigma^2$); and -- this is the special
part -- the \emph{shape stays normal}. Shifting and stretching a bell gives another
bell. For almost any other distribution this fails: stretch an exponential and you
get another exponential only if you do not shift it; the family is far more
fragile. The normal family is closed under \emph{every} non-degenerate linear map.

\subsection{Sums of independent normals are normal}
\label{ch12:sums}

The deeper and more surprising closure property:

\begin{theorem}[Additive closure]
\label{ch12:thm-sum}
If $X \sim \Normal(\mu_1, \sigma_1^2)$ and $Y \sim \Normal(\mu_2, \sigma_2^2)$ are
\emph{independent}, then
\[
  X + Y \sim \Normal\!\big(\mu_1 + \mu_2,\; \sigma_1^2 + \sigma_2^2\big).
\]
\end{theorem}

The means add (always true for sums); the variances add (always true for
\emph{independent} sums); and again the shape stays normal -- which is special. The
sum of two independent normals is not some new lumpy thing; it is just a wider
bell.

\begin{whybox}
Why do the variances add but not the standard deviations? Variance is the natural
additive quantity for independent sums:
\[
  \Var(X+Y) = \Var(X) + \Var(Y) \quad (\text{independence}),
\]
because the cross-term $2\Cov(X,Y)$ vanishes when $X$ and $Y$ are independent.
Standard deviations therefore add ``in quadrature'':
$\sigma_{X+Y} = \sqrt{\sigma_1^2 + \sigma_2^2}$, like the hypotenuse of a right
triangle with legs $\sigma_1$ and $\sigma_2$. Two equal-sized independent
errors of size $\sigma$ do \emph{not} combine into an error of size $2\sigma$;
they combine into $\sqrt{2}\,\sigma \approx 1.41\sigma$. Random errors partly
cancel.

That the result is again \emph{normal} is the genuinely deep fact. The cleanest
proof uses the \emph{moment generating function} (an algebraic transform
$M_X(t) = \E[e^{tX}]$ that packages all the moments $\E[X^k]$ of a distribution
into a single generating function, explored fully in Chapter 15): for
$X\sim\Normal(\mu,\sigma^2)$,
\[
  M_X(t) = \E[e^{tX}] = \exp\!\Big(\mu t + \tfrac{1}{2}\sigma^2 t^2\Big).
\]
For independent $X,Y$ the MGF of the sum is the product of the MGFs:
\[
  M_{X+Y}(t) = M_X(t)\,M_Y(t)
  = \exp\!\Big((\mu_1+\mu_2)t + \tfrac{1}{2}(\sigma_1^2+\sigma_2^2)t^2\Big),
\]
which is precisely the MGF of a $\Normal(\mu_1+\mu_2,\sigma_1^2+\sigma_2^2)$. Since
the MGF determines the distribution, the sum \emph{is} that normal. Adding the
quadratic-in-$t$ exponents is what adds the variances; the form
``$\exp(\text{linear} + \text{quadratic})$'' is closed under multiplication, and
that algebraic closure is the secret behind the statistical closure.
\end{whybox}

\begin{figure}[H]
\centering
\begin{tikzpicture}
\begin{axis}[
  width=12cm, height=6cm,
  domain=-6:10, samples=180,
  xlabel={$x$}, ylabel={density},
  ymin=0, ymax=0.45,
  axis lines=middle, enlargelimits=false,
  title style={font=\small},
  title={Additive closure: sum of independent normals},
  legend pos=north east, legend cell align=left,
]
\addplot[distblue, very thick] {1/(1*sqrt(2*pi))*exp(-(x-0)^2/(2*1^2))};
\addlegendentry{$X\sim\Normal(0,1)$}
\addplot[targetgreen, very thick, dashed] {1/(1.5*sqrt(2*pi))*exp(-(x-2)^2/(2*1.5^2))};
\addlegendentry{$Y\sim\Normal(2,2.25)$}
\addplot[densityred, very thick, dotted]
  {1/(sqrt(1+2.25)*sqrt(2*pi))*exp(-(x-2)^2/(2*(1+2.25)))};
\addlegendentry{$X+Y\sim\Normal(2,3.25)$}
\end{axis}
\end{tikzpicture}
\caption{Adding two independent normals gives a third normal: means add
($0+2=2$), variances add ($1 + 2.25 = 3.25$), and the sum is wider than either
piece but still a clean bell.}
\label{ch12:fig-sum}
\end{figure}

\begin{example}[Stacking tolerances]
\label{ch12:ex-stack}
A factory stacks two metal plates. Plate A has thickness $\Normal(10, 0.2^2)$ mm
and plate B has thickness $\Normal(15, 0.3^2)$ mm, independently. The total
thickness is
\[
  T = A + B \sim \Normal\big(10+15,\; 0.2^2 + 0.3^2\big) = \Normal(25,\; 0.13).
\]
So the stack averages $25$ mm with standard deviation $\sqrt{0.13}\approx 0.36$ mm.
What is the chance the stack exceeds $25.5$ mm?
\[
  \P(T > 25.5) = \P\!\left(Z > \frac{25.5 - 25}{0.36}\right)
              = \P(Z > 1.39) \approx 1 - 0.9177 = 0.0823.
\]
About $8\%$. Note the combined sd $0.36$ is less than $0.2 + 0.3 = 0.5$: errors
partly cancel, just as the in-quadrature rule promised.
\end{example}

\begin{corollary}[Averages of normals]
\label{ch12:cor-average}
If $X_1,\dots,X_n$ are independent $\Normal(\mu,\sigma^2)$, then their sum is
$\Normal(n\mu, n\sigma^2)$ and their average is
\[
  \bar X = \frac{1}{n}\sum_{i=1}^n X_i \sim \Normal\!\left(\mu, \frac{\sigma^2}{n}\right).
\]
\end{corollary}

This little corollary is the engine of all of statistics. Averaging $n$ normal
measurements gives a normal answer centered at the true $\mu$ with the spread
shrunk by a factor $\sqrt{n}$. Average four measurements and you halve the error;
average a hundred and you cut it tenfold. The astronomer from
Section~\ref{ch12:errors} was right to take many readings and average them.

\begin{walkbox}
\textbf{Mental check: independent error cancellation.}
Suppose you measure a physical length with a digital caliper that has normal
error with standard deviation $\sigma = 10\,\mu\text{m}$. You take two
independent readings and average them: $\bar X = (X_1 + X_2)/2$. What is the
standard deviation of your average? Is it $5\,\mu\text{m}$,
$7.07\,\mu\text{m}$, or $10\,\mu\text{m}$? (Answer: By
Corollary~\ref{ch12:cor-average}, the variance of the average is $\sigma^2/2 =
100/2 = 50$, so the standard deviation is $\sqrt{50} = 10/\sqrt{2} \approx
7.07\,\mu\text{m}$. It does not drop to $5\,\mu\text{m}$ until you average
$n=4$ readings! Precision improves with the square root of sample size,
$\sqrt{n}$, not linearly.)
\end{walkbox}

\section{Why Is the Normal Everywhere? A Preview of the CLT}
\label{ch12:everywhere}

We can now address the mystery in the chapter's introduction. Why does this one
curve describe heights, errors, noise, test scores, and the daily wiggles of a
stock index -- things with no obvious connection to each other?

The answer is that all of those quantities are \emph{sums of many small,
independent effects}. A person's height is the cumulative result of thousands of
genes and environmental nudges, each contributing a tiny push up or down. A
measurement error is the sum of many small disturbances -- air, vibration, the
observer's nerves. The key result, which Chapter~21 makes rigorous, is:

\begin{aestheticbox}
\textbf{The Central Limit Theorem (preview).} When you add up many independent
random contributions, no single one of them dominating, the distribution of the
sum -- after centering and rescaling -- looks like a normal distribution,
\emph{almost regardless of what the individual contributions look like}. The
ingredients can be coin flips, dice, skewed payouts, anything with a finite
variance. Their accumulated sum forgets its origins and converges to the bell.
\end{aestheticbox}

This is the real reason the normal is ``normal.'' It is not that nature has a
fondness for $e^{-x^2/2}$. It is that nature is full of \emph{accumulation}, and
accumulation of many small independent pieces forces the bell to appear. The
normal is the universal shape of \emph{averages and aggregates}.

We can already see the convergence happening with the humblest of building blocks:
sums of coin flips. Add up $n$ independent $\pm 1$ steps and the histogram of the
result, once $n$ is even modestly large, is visually indistinguishable from a
normal curve.

\begin{figure}[H]
\centering
\begin{tikzpicture}
\begin{axis}[
  width=13cm, height=6cm,
  ybar interval=0.95,
  xlabel={value of the sum (standardized)},
  ylabel={relative frequency},
  ymin=0, ymax=0.45,
  xmin=-3.5, xmax=3.5,
  axis lines=left,
  ytick={0.2, 0.4},
  xtick={-3,-2,-1,0,1,2,3},
  title style={font=\small},
  title={Convergence of $12$ coin flips to the standard normal bell}
]
% Binomial(n=12, p=1/2) standardized: mean 6, sd sqrt(3)~1.732,
% bars rescaled to density by dividing pmf by the z-bin width 1/1.732.
\addplot[fill=histfill, draw=distblue!70!black] coordinates {
  (-3.753,0.0004) (-3.175,0.0051) (-2.598,0.0279) (-2.021,0.0930)
  (-1.443,0.2093) (-0.866,0.3349) (-0.289,0.3907) (0.289,0.3349)
  (0.866,0.2093) (1.443,0.0930) (2.021,0.0279) (2.598,0.0051)
  (3.175,0.0004) (3.753,0.0000)
};
\addplot[densityred, very thick, samples=120, domain=-3.5:3.5]
  {1/sqrt(2*pi)*exp(-x^2/2)};
\end{axis}
\end{tikzpicture}
\caption{The standardized distribution of a sum of $12$ fair coin steps (bars,
each $\pm1$) against the standard normal density (curve). With only twelve flips
the bell is already an excellent fit -- this is de Moivre's original discovery and
a glimpse of the Central Limit Theorem.}
\label{ch12:fig-clt}
\end{figure}

\begin{figure}[H]
\centering
\begin{tikzpicture}[scale=0.9]
% Funnel at the top
\draw[thick, black!70] (-1.0,4.6) -- (-0.25,3.9) -- (-0.25,3.4);
\draw[thick, black!70] (1.0,4.6) -- (0.25,3.9) -- (0.25,3.4);
\node[above, font=\footnotesize\bfseries, black!80] at (0,4.6) {Bead Drop};

% Triangular peg array
\foreach \row in {0,...,5} {
  \foreach \col in {0,...,\row} {
    \fill[black!75] ({(\col - \row/2)*0.95}, {3.0 - \row*0.55}) circle (2.2pt);
  }
}

% Path sample (trajectory showing independent left/right choices)
\draw[thick, trapred, -Stealth] (0,3.4) -- (0,3.0);
\draw[thick, trapred, -Stealth] (0,3.0) -- (0.475,2.45);
\draw[thick, trapred, -Stealth] (0.475,2.45) -- (0.0,1.9);
\draw[thick, trapred, -Stealth] (0.0,1.9) -- (0.475,1.35);
\draw[thick, trapred, -Stealth] (0.475,1.35) -- (0.0,0.8);
\draw[thick, trapred, -Stealth] (0.0,0.8) -- (0.475,0.25);
\node[right, font=\tiny, trapred!90!black] at (0.5, 1.35) {$\pm 1$ step at each pin};

% Collection bins at the base
\foreach \b in {-3,...,4} {
  \draw[thick, black!60] ({\b*0.95 - 0.475}, -1.8) -- ({\b*0.95 - 0.475}, 0.1);
}
\draw[thick, black!80] (-3.325, -1.8) -- (3.325, -1.8);

% Accumulated histogram bars (Binomial n=6)
\fill[histfill, draw=distblue!80!black, rounded corners=1pt] (-3.325,-1.8) rectangle (-2.375, -1.74);
\fill[histfill, draw=distblue!80!black, rounded corners=1pt] (-2.375,-1.8) rectangle (-1.425, -1.45);
\fill[histfill, draw=distblue!80!black, rounded corners=1pt] (-1.425,-1.8) rectangle (-0.475, -0.75);
\fill[histfill, draw=distblue!80!black, rounded corners=1pt] (-0.475,-1.8) rectangle (0.475, -0.32);
\fill[histfill, draw=distblue!80!black, rounded corners=1pt] (0.475,-1.8) rectangle (1.425, -0.75);
\fill[histfill, draw=distblue!80!black, rounded corners=1pt] (1.425,-1.8) rectangle (2.375, -1.45);
\fill[histfill, draw=distblue!80!black, rounded corners=1pt] (2.375,-1.8) rectangle (3.325, -1.74);

% Overlaid normal bell curve
\draw[very thick, densityred, domain=-3.3:3.3, samples=60]
  plot (\x, {-1.8 + 1.55*exp(-\x*\x/1.9)});

\node[font=\footnotesize, densityred!90!black, right] at (1.6, -0.2) {Continuous Bell Limit $\Normal(\mu,\sigma^2)$};
\node[font=\scriptsize, black!70, below] at (0,-1.85) {Bins collect the sum of independent left/right deflections};
\end{tikzpicture}
\caption{Galton's Quincunx (the bean machine). Beads drop through a triangular grid of pegs, randomly deflecting left or right with probability $1/2$ at each stage. By the time a bead reaches the bottom, its horizontal position is the sum of independent binary choices. The accumulated heap forms the unmistakable bell of the normal distribution.}
\label{ch12:fig-galton}
\end{figure}

\begin{connectionbox}
\textbf{Making probability visible: Galton's machine.}
In 1873, Francis Galton built this physical apparatus -- which he named the
\emph{quincunx} after the geometric pattern of five dots on a die -- to make the
abstract mathematics of the bell curve visible to the human eye. As thousands of
tiny lead pellets drop through the funnel and cascade past the triangular
pegboard, each ball acts out a sequence of independent binary choices, deflecting
left or right with equal probability. When the shot settles into the vertical
slots at the base, the random scatter unfailingly sculpts the smooth, symmetric
dome of the normal density. Galton was captivated by this emergence of order from
chaos: no single pellet's path can be predicted, yet the collective heap is as
reproducible as a law of gravity.
\end{connectionbox}

\section{A Paradox and a Caution}
\label{ch12:paradox}

The normal distribution is so useful that it is over-used, and its thin tails are
the most common place where the model and reality part ways.

\begin{paradoxbox}
\textbf{The ``impossible'' market crash.} Suppose daily returns of a stock index
are modeled as normal, and the historical daily standard deviation is about $1\%$.
On October 19, 1987 (``Black Monday''), the index fell about $20\%$ in a single
day -- a $20$-sigma event.

How rare is a $20$-sigma event under a normal model? The tail probability
$\P(Z < -20)$ is on the order of $10^{-89}$. To put that in perspective: if the
market had been open every day since the Big Bang, roughly $5\times 10^{12}$ days,
the expected number of such days would still be about $10^{-76}$ -- effectively
zero, short of one by some seventy orders of magnitude. Under the normal model,
Black Monday simply cannot happen.

But it \emph{did} happen -- and so did several other ``impossible'' days. The
resolution is not that probability is broken. It is that the \emph{model is
wrong}: real financial returns are \emph{not} normal. They have \emph{fat tails}
-- extreme moves are far more common than a bell curve predicts. The normal's
$e^{-z^2/2}$ decay is extremely rapid. While this property makes the distribution
mathematically tractable, it makes it an unrealistic model for phenomena where rare,
severe events dominate. The lesson: the normal is the appropriate model for
\emph{averages of many small independent effects}, whereas market shocks are
rare, correlated, and dominant.
\end{paradoxbox}

\begin{trapbox}
``Normal'' does not mean ``default.'' Before assuming a quantity is normal, ask
the CLT question: \emph{is this an average or sum of many small independent
pieces?} Heights, measurement errors, and sample means: yes, plausibly normal.
Incomes (a few people earn enormous amounts): no, heavily skewed. Waiting times:
no, often exponential. City sizes and word frequencies: no, power-law. Reaching
for the bell out of habit is one of the most common modeling mistakes in applied
statistics.
\end{trapbox}

\begin{applicationbox}
\textbf{Where the normal genuinely rules in machine learning.} Even with the
caution above, the Gaussian is standard in modern statistics and machine
learning for good reasons. In linear regression, minimizing squared prediction
errors is mathematically identical to \emph{maximum-likelihood estimation} (the
principle of picking parameters that make the observed data most probable) under
the assumption of additive normal noise. \emph{Gaussian mixture models} cluster
complex datasets by representing them as weighted combinations of several
overlapping bell curves. \emph{Gaussian processes} extend the bell curve to
function spaces, placing a prior distribution over entire smooth curves for
Bayesian optimization. And in deep learning, the \emph{reparameterization trick}
used in variational autoencoders (generative neural networks that learn
compressed latent representations of complex data) relies on writing a random
latent variable as $X = \mu + \sigma Z$ with fixed standard noise $Z \sim \Normal(0,1)$.
This is the standardization identity of Theorem~\ref{ch12:thm-standardize} run in
reverse, allowing calculus gradients to flow smoothly through the parameters
$\mu$ and $\sigma$ without getting blocked by the random sampling step. In every
one of these applications, the linear-closure and additive-closure properties of
this chapter are what make the calculations possible.
\end{applicationbox}

\section{The Bivariate Normal, Briefly}
\label{ch12:bivariate}

Normals generalize to several dimensions, and the two-dimensional case is worth a
glimpse because it reveals what the curve looks like when two normal quantities
travel together. If $X$ and $Y$ are jointly normal, their density is a bell over
the plane, and its level sets -- the curves of equal height -- are \emph{ellipses}.
When $X$ and $Y$ are independent the ellipses are axis-aligned; when they are
correlated the ellipses tilt, leaning along the direction of the correlation.

\begin{figure}[H]
\centering
\begin{tikzpicture}
\begin{axis}[
  width=8.5cm, height=8.5cm,
  xlabel={$x$}, ylabel={$y$},
  xmin=-3.2, xmax=3.2, ymin=-3.2, ymax=3.2,
  xtick=\empty, ytick=\empty,
  axis lines=middle, enlargelimits=false,
  title style={font=\small},
  title={Level sets of a correlated bivariate normal},
]
% Tilted ellipses (correlated normal contours): an ellipse with semi-axes
% (1.6, 0.7) rotated by 45 degrees, drawn at three sizes.
\foreach \s in {0.7,1.4,2.1} {
  \addplot[distblue, thick, domain=0:360, samples=80]
    ({\s*(1.6*cos(x)*0.7071 - 0.7*sin(x)*0.7071)},
     {\s*(1.6*cos(x)*0.7071 + 0.7*sin(x)*0.7071)});
}
\addplot[densityred, dashed, thick] coordinates {(-3,-3)(3,3)};
\node[densityred] at (axis cs:2.2,1.5) {\footnotesize tilt};
\end{axis}
\end{tikzpicture}
\caption{Contours of a bivariate normal with positive correlation. The tilted,
elongated ellipses show that large $x$ tends to come with large $y$. Independent
normals would give circles or axis-aligned ellipses.}
\label{ch12:fig-bivariate}
\end{figure}

\begin{connectionbox}
The full multivariate normal is governed by a mean \emph{vector} $\boldsymbol\mu$
(the list of individual center coordinates) and a \emph{covariance matrix} $\Sigma$
(a square grid recording how much each pair of variables fluctuates together),
with density proportional to
$\exp\!\big(-\tfrac12 (\mathbf{x}-\boldsymbol\mu)^\top \Sigma^{-1}(\mathbf{x}-\boldsymbol\mu)\big)$.
The quadratic form in the exponent is the multidimensional cousin of
$(x-\mu)^2/\sigma^2$, and the ellipses of Figure~\ref{ch12:fig-bivariate} are its
level sets -- curves of equal density whose tilts and widths are dictated by the
eigenvectors (perpendicular axes of stretch) and eigenvalues (stretch factors) of
$\Sigma$. This is where probability and linear algebra meet: rotating the
coordinate axes along the eigenvectors uncouples the correlated variables into
independent standard normals. We develop this machinery fully in the coming
chapters on joint distributions and covariance.
\end{connectionbox}

\section{What to Carry Forward}
\label{ch12:summary}

\begin{recoverybox}
The essential picture:
\begin{itemize}
  \item The standard normal density is $\varphi(z) = \frac{1}{\sqrt{2\pi}}e^{-z^2/2}$.
  The $e^{-z^2/2}$ is the bell; the $\sqrt{2\pi}$ is just the area-normalizer, and
  it is $\sqrt{2\pi}$ because of the polar-coordinates Gaussian integral.
  \item The general normal $\Normal(\mu,\sigma^2)$ is the same curve slid to $\mu$
  and stretched by $\sigma$. Standardizing, $Z = (X-\mu)/\sigma$, turns every
  normal question into one lookup of $\Phi$.
  \item $68$--$95$--$99.7$: the chance of being within one, two, or three standard
  deviations of the mean, for \emph{any} normal.
  \item Linear maps of a normal stay normal; sums of independent normals stay
  normal, with means adding and variances adding. These closure properties are
  what make the normal so easy to compute with.
  \item The normal is everywhere because of the Central Limit Theorem: it is the
  universal shape of sums and averages of many small independent effects. But its
  thin tails make it a poor model for rare, extreme, correlated shocks.
\end{itemize}
\end{recoverybox}

In Section~\ref{ch12:bivariate}, we saw two random variables traveling together across the plane.
Up to now, each random variable has lived on its own number line.
The next chapter examines what happens when multiple random variables interact:
their joint probabilities, dependencies, and conditional distributions.

\section*{Exercises}
\addcontentsline{toc}{section}{Exercises}

\begin{exercise}[Reading the curve]
\label{ch12:exr-read}
Let $Z \sim \Normal(0,1)$. Using the table values in Section~\ref{ch12:table} and
the symmetry/complement rules, compute:
(a) $\P(Z > 2)$; (b) $\P(Z < -0.5)$; (c) $\P(-1 \le Z \le 2)$;
(d) $\P(|Z| > 1)$.
(e) Without computing, which probability is larger: $\P(0 \le Z \le 1)$ or $\P(1 \le Z \le 2)$? Explain using the geometric shape of the density curve.
\end{exercise}

\begin{exercise}[Reading parameters from geometry]
\label{ch12:exr-inflection}
An unlabeled plot of a normal density curve has its highest peak at $x = 50$,
and the curvature changes from concave down (the central dome) to concave up (the
flaring tails) at $x = 42$ and $x = 58$.
(a) What are the mean $\mu$ and standard deviation $\sigma$ of this distribution?
(b) What is the maximum peak height of the density curve?
(c) Using the empirical rule, between what two values does approximately $95\%$
of the probability mass lie?
\end{exercise}

\begin{exercise}[Standardizing]
\label{ch12:exr-standardize}
Adult systolic blood pressure is modeled as $X \sim \Normal(120, 12^2)$ mmHg.
(a) What fraction of adults have pressure above $140$ (a common ``high'' cutoff)?
(b) Below $100$? (c) Between $108$ and $132$? Relate part (c) to the empirical
rule.
\end{exercise}

\begin{exercise}[Inverse lookups]
\label{ch12:exr-inverse}
For $Z \sim \Normal(0,1)$, find the cutoffs (using the table or known critical
values): (a) the value $c$ with $\P(Z > c) = 0.025$; (b) the value $c$ with
$\P(Z \le c) = 0.90$; (c) the symmetric interval $[-c, c]$ containing exactly
$99\%$ of the mass.
\end{exercise}

\begin{exercise}[Adding normals]
\label{ch12:exr-sum}
Let $X \sim \Normal(3, 4)$ and $Y \sim \Normal(-1, 5)$ be independent.
Find the distribution of (a) $2X - 1$; (b) $X + Y$; (c) $X - Y$;
(d) $\tfrac{1}{2}(X + Y)$. In each case give the mean and variance and state why
the result is still normal.
\end{exercise}

\begin{exercise}[The shrinking error]
\label{ch12:exr-average}
A single thermometer reading is $\Normal(\mu, 0.4^2)$ degrees, where $\mu$ is the
true temperature. (a) How many independent readings must you average so that the
standard deviation of the average is below $0.1$ degrees? (b) With that many
readings, what is the chance the average is within $0.1$ of the truth? Use
Corollary~\ref{ch12:cor-average}.
\end{exercise}

\begin{exercise}[Why $\sqrt{2\pi}$, your turn]
\label{ch12:exr-gaussint}
Show that $\int_{-\infty}^{\infty} e^{-a x^2}\,\dd x = \sqrt{\pi/a}$ for $a > 0$,
by the polar-coordinates square trick of Theorem~\ref{ch12:thm-gaussint}. Then set
$a = 1/(2\sigma^2)$ to confirm that
\[
  \frac{1}{\sigma\sqrt{2\pi}}\int_{-\infty}^{\infty} e^{-(x-\mu)^2/(2\sigma^2)}\,\dd x = 1,
\]
i.e.\ that the general normal density really integrates to $1$.
\end{exercise}

\begin{exercise}[A proof: oddness gives the mean]
\label{ch12:exr-mean}
Prove carefully that $\E[Z] = 0$ for $Z \sim \Normal(0,1)$ by splitting
$\int_{-\infty}^{\infty} z\,\varphi(z)\,\dd z$ into $\int_{-\infty}^{0}$ and
$\int_{0}^{\infty}$ and using the substitution $z \mapsto -z$ on one piece. Where
exactly does the symmetry (evenness of $\varphi$) get used?
\end{exercise}

\begin{exercise}[Variance by parts]
\label{ch12:exr-var}
Reproduce the integration-by-parts computation of $\Var(Z) = 1$ from
Proposition~\ref{ch12:prop-moments}, filling in the missing step: explain in your
own words why the boundary term $\big[-z\,e^{-z^2/2}\big]_{-\infty}^{\infty}$ is
zero. (Hint: compare the growth of $z$ to the decay of $e^{-z^2/2}$.)
\end{exercise}

\begin{exercise}[The fat-tail check]
\label{ch12:exr-tails}
A risk model assumes daily losses are $\Normal(0, (1\%)^2)$. (a) Under this model,
what is the probability of a daily loss exceeding $5\%$? (b) Exceeding $10\%$?
(c) Empirically, large markets see $5\%$+ down days every few years. Roughly how
many years between $5\%$ down days does the normal model predict (assume $250$
trading days per year)? Comment on what this says about the normal as a model for
financial risk, connecting to the paradox box in Section~\ref{ch12:paradox}.
\end{exercise}

\begin{exercise}[Conceptual: shape versus parameters]
\label{ch12:exr-concept}
Explain, in a few sentences each and without heavy algebra: (a) why the peak of a
$\Normal(\mu,\sigma^2)$ density gets \emph{lower} as $\sigma$ grows; (b) why
standard deviations of independent normals add ``in quadrature'' rather than
straight; (c) why a single function $\Phi$ suffices to answer questions about
every normal distribution, regardless of $\mu$ and $\sigma$.
\end{exercise}

\begin{exercise}[The maximum-entropy view, challenge]
\label{ch12:exr-maxent}
It is a deep fact that among all distributions on $\R$ with a given mean $\mu$ and
variance $\sigma^2$, the normal is the one with the \emph{largest entropy} -- the
``most spread out'' or ``least committal'' such distribution. Without proving it,
explain in words why this is an appealing justification for using the normal as a
default model of a quantity when you know only its mean and variance, and connect
it to the error-curve story of Section~\ref{ch12:errors}. (We return to entropy in
the information-theory chapter.)
\end{exercise}
